% ============================================================
% APPENDIX F -- Verification: the code that asserts the results
% ============================================================

\label{app:verification}

\begin{ataglance}[How to read this appendix]
Every numbered result in this thesis that can be checked by computation is
checked by computation, and this appendix is the record.

\S\ref{sec:F_why} says what the suite does and, more importantly, what it does
not do: it verifies mathematics, not medicine, and passing it is not evidence
about any patient. \S\ref{sec:F_arch} describes the reference implementation
and the naming rule that ties one test to one numbered statement.
\S\ref{sec:F_run} says how to run it. \S\ref{sec:F_results} prints the outcome
of the run that produced this copy of the thesis. \S\ref{sec:F_found} records
the two errors the suite caught, because a verification appendix that reports
only successes is a marketing document. \S\ref{sec:F_limits} sets out what
remains unverified and why.
\end{ataglance}

\section{What the Suite Is For, and What It Cannot Do}
\label{sec:F_why}

A proof establishes that a statement follows from its premises. It does not
establish that the statement was transcribed correctly into the thesis, that
the arithmetic in a worked example is right, or that the implementation someone
later writes computes the object the proof was about. Those are different
failures, they are common, and they are the ones a test suite catches.

The suite here does four things. It checks that each stated property holds for
the implemented function across a grid of inputs, which catches a
transcription error in a formula. It checks that each worked example in the
text returns the number the text prints, which catches arithmetic slips and,
more usefully, catches them again whenever a parameter is later edited. It
checks the internal consistency of the running cases, so that Mrs~A's realised
capability is the same figure in Chapter~\ref{ch:theory} and in
Chapter~\ref{app:crf}. And it checks the qualitative claims the prose makes
about the mathematics, such as that lowering the substitution parameter
punishes imbalance harder, which are the claims most easily asserted and least
often verified.

What it does not do needs stating plainly, because verification suites attract
more credit than they deserve.

\begin{itemize}
\item \textbf{It is not a proof.} A property that holds on a finite grid of
inputs may fail elsewhere. The proofs are in
Appendix~\ref{app:mathematical_foundations} and Chapter~\ref{app:crf}; the tests check
that the implementation agrees with them, which is a weaker and different
claim.

\item \textbf{It is not empirical validation.} No participant data are used.
Nothing here says whether the three domains exist, whether the indicators
measure them, or whether the construct predicts anything. Those questions
belong to Chapters~\ref{ch:results1} and~\ref{ch:results2} and are answered by
data, not by assertions in Python.

\item \textbf{It cannot check the axioms.} Whether capacity boundedness is the
right commitment, or whether $\varepsilon$ should be $0$ or $2$, are questions
no test can settle. The suite verifies that the mathematics does what the
axioms say; it is silent on whether the axioms are correct.

\item \textbf{Passing is not a quality claim.} A suite tests what its author
thought to test. Its value lies in what it caught, which
\S\ref{sec:F_found} reports, and in the fact that it will catch the same class
of error again when the thesis is next edited.
\end{itemize}

The honest summary is that this appendix raises confidence in the
\emph{bookkeeping} of a long mathematical argument. That is worth having. It is
not worth more than it is.

\section{Architecture and the Naming Rule}
\label{sec:F_arch}

The code lives in \texttt{verification/} in the thesis archive and has two
parts. The package \texttt{hpv} is a reference implementation: one module per
mathematical layer, with every function carrying the number of the result it
implements. The directory \texttt{tests/} holds the assertions.

\begin{table}[htbp]\centering
\caption{The reference implementation. Each module implements the objects of
one mathematical layer of the thesis.}
\label{tab:F_modules}
\footnotesize
\begin{tabular}{p{3.4cm} p{2.4cm} p{7.6cm}}
\toprule
\textbf{Module} & \textbf{Source} & \textbf{Objects implemented} \\
\midrule
\texttt{hpv/crf.py} & App.~\ref{app:mathematical_foundations}, Ch.~\ref{app:crf} &
Domain profile, CES contextual mean and its geometric limit, the additive and
relational comparators, deprivation counting, gradients and elasticities,
cross-partials, the gap decomposition, the Atkinson equally distributed
equivalent and its social weights \\
\addlinespace
\texttt{hpv/dynamics.py} & App.~\ref{app:dynamical} &
Michaelis--Menten and Hill terms, the drift of the deterministic model, its
closed-form equilibrium, the Lipschitz constant, fixed-step RK4,
Euler--Maruyama, and a conservative finite-volume Fokker--Planck step \\
\addlinespace
\texttt{hpv/haemophilia.py} & App.~\ref{app:hemo_dynamics} &
Neutralisation fraction, the factor steady state, the inhibitor threshold, the
linear reference model, the ITI equilibrium and its critical suppression, the
gene-therapy equilibrium \\
\addlinespace
\texttt{hpv/economics.py} & App.~\ref{app:a-economics}, \ref{app:computation} &
Net health benefit, incremental net monetary benefit, displaced health, budget
impact, discounting, the ICER, and the DALY identity \\
\addlinespace
\texttt{hpv/prognostic.py} & Ch.~\ref{ch:methods}, \ref{ch:results2} &
Hazard-ratio to absolute-risk conversion, cumulative incidence under a
competing risk against the Kaplan--Meier complement, decision-curve net
benefit, and classical-error attenuation \\
\bottomrule
\end{tabular}
\end{table}

The naming rule is the part that makes the suite usable rather than
decorative. Every test function is named for the statement it asserts:
\texttt{test\_proposition\_inhibitor\_threshold},
\texttt{test\_theorem\_gap\_decomposition\_is\_exact},
\texttt{test\_worked\_example\_mrs\_a\_absolute\_risk}. A reader who doubts
Proposition~\ref{prop:C_threshold} can find the test that asserts it by name,
read the assertion, and change the numbers. Every test carries a docstring
naming the result and saying in one sentence what would be wrong if it failed.

This is also why the pointer beneath each proposition in the appendices reads
\emph{Verified numerically by} and then a test name. In the previous version of
this thesis those pointers named tests in a package called \texttt{volos} that
did not exist. Announcing a verification engine one has not written is worse
than writing nothing, and the pointers now name tests that run.

\section{Running the Suite}
\label{sec:F_run}

The suite needs Python 3.10 or later, \texttt{pytest} and \texttt{numpy}, and
nothing else. From the \texttt{verification/} directory:

\begin{lstlisting}[style=hpcode,language=bash]
pip install pytest numpy
python3 -m pytest tests/ -q          # run everything
python3 -m pytest tests/ -q -k threshold   # run one result
python3 run_verification.py          # run everything and write the table below
\end{lstlisting}

\noindent
The runner writes two files. \texttt{verification\_results.tex} is the table in
\S\ref{sec:F_results}, regenerated on every build, so the table cannot drift
away from the code that produced it. \texttt{verification\_results.json} is the
same record in machine-readable form, with the timestamp, the Python version
and the platform, so a reader can tell whether the printed table came from the
code they are holding.

Stochastic tests fix their seeds. The Monte Carlo comparisons use tolerances
set from the standard error of the estimator rather than tuned until they
passed, and the seeds are declared in the source so a failure can be
reproduced.

\section{Results of the Run That Produced This Document}
\label{sec:F_results}

% Generated by verification/run_verification.py. Do not edit by hand.
\begin{table}[htbp]\centering
\caption{Result of running the verification suite. Each test asserts a numbered statement from the chapter or appendix named. The suite is re-run and this table regenerated on every build of the thesis.}
\label{tab:verification}
\footnotesize
\begin{tabular}{p{4.6cm} p{1.6cm} p{5.4cm} r r}
\toprule
\textbf{Test module} & \textbf{Covers} & \textbf{Subject} & \textbf{Passed} & \textbf{Failed} \\
\midrule
\texttt{test\_axioms.py} & A & Axioms, aggregation and the rival ontology & 39 & 0 \\
\texttt{test\_crf.py} & E & Capability Realisation Function & 19 & 0 \\
\texttt{test\_dynamics.py} & B & Dynamical layers: ODE, SDE, Fokker--Planck & 14 & 0 \\
\texttt{test\_haemophilia.py} & C & Factor kinetics, ITI and gene therapy & 11 & 0 \\
\texttt{test\_economics.py} & A, D & Net benefit, budget impact and DALYs & 10 & 0 \\
\texttt{test\_prognostic.py} & 3, 6 & Survival, competing risks and net benefit & 16 & 0 \\
\texttt{test\_measurement.py} & 5 & Reliability, validity and factor scores & 10 & 0 \\
\texttt{test\_chapter\_examples.py} & 1, 7--9 & Worked examples in the chapters & 7 & 0 \\
\midrule
\textbf{Total} & & & \textbf{126} & \textbf{0} \\
\bottomrule
\end{tabular}
\end{table}

\noindent\emph{Run 2026-09-03 11:58 UTC on Python 3.11.15, Linux-6.18.44-fc-v24-x86\_64-with-glibc2.39, in 30.1\,s.}


\noindent
A failing row would appear here, in this table, in the submitted thesis. That
is the intention. A verification appendix whose table is generated by the same
command that builds the document cannot quietly report a state of affairs that
has ceased to be true.

\section{What the Suite Caught}
\label{sec:F_found}

Two errors surfaced on the first runs. Both are recorded because the value of a
test suite is measured in what it finds, and because both are instructive about
what kinds of checking catch what kinds of mistake.

\subsection{A worked example that proved nothing}

The compensation example in \S\ref{app:a-additivity} was written with Mr~B's
structural freedom falling from $0.20$ to $0.05$ while immune tolerance
induction raised his biomedical potential from $0.30$ to $0.60$. The additive
index scores the change as an improvement, which was the point. The realisation
function was supposed to score it as a deterioration, and does not: both
profiles return exactly $0.112$, because
$0.60\sqrt{0.70\times0.05}$ and $0.30\sqrt{0.70\times0.20}$ are the same
number. The example demonstrated a difference between the two rules that, at
those particular numbers, did not exist.

The parameter was changed to $0.04$, where the realisation function returns
$0.100$ against $0.112$ and the two rules disagree about the sign of the
change. The point survives; the arithmetic that was supposed to carry it now
does.

The general lesson is that a hand-worked example is a piece of evidence, and
like any piece of evidence it can be checked or taken on trust. This one had
been read several times by a human and the coincidence was not noticed.

\subsection{A sign error that mass conservation could not see}

The finite-volume Fokker--Planck step in \texttt{hpv/dynamics.py} was written
with the advective flux entering the assembled operator with the wrong sign, so
a positive drift moved the population density towards lower biomarker values
instead of higher ones.

The interesting part is which test failed. The mass-conservation test of
Proposition~\ref{prop:B_mass} passed throughout, both before and after the
correction, because a sign error in the flux still leaves the operator's column
sums at zero: probability was conserved perfectly while flowing the wrong way.
Only the test asserting that drift moves the mean and diffusion widens the
variance, which is the qualitative claim the prose in \S\ref{sec:B_fp} makes,
caught it.

Conservation properties are the natural thing to test in a discretised
transport problem, and they are insensitive to direction. A suite built only
from the invariants a numerical analyst would think to check would have shipped
this. The test that caught it was written to verify a sentence in the
manuscript rather than a property of the scheme, which is an argument for
writing tests from the prose as well as from the mathematics.

\section{What Remains Unverified}
\label{sec:F_limits}

Three classes of claim in this thesis are not covered, and it is better to name
them than to let the table above imply a coverage it does not have.

\textbf{Everything that needs data.} The measurement model, the factor
structure, the invariance tests, the survival models and every performance
statistic in Chapters~\ref{ch:results1} and~\ref{ch:results2} require the NHATS
extract, which is not part of this archive. The R scripts under
\texttt{analysis/} validate their inputs and record session information but
cannot reproduce a numerical result until a documented extract and variable
mapping are supplied. What the Python suite covers on the prognostic side is
the arithmetic of hazard conversion, competing risks, net benefit and
attenuation, which is the part that can be checked without a cohort.

\textbf{The proofs themselves.} No theorem prover is used. The tests check
agreement between the implementation and the stated result on sampled inputs;
they cannot establish that a proof is valid, and a reader who doubts a proof
should read the proof.

\textbf{The normative content.} The axioms, the choice of $\alpha$, the range
of $\rho$, the value of $\varepsilon$ and the opportunity-cost threshold
$\lambda$ are judgements. The suite verifies what follows from a given choice.
It has nothing to say about whether the choice is defensible, and no amount of
green output should be read as though it did.
